\documentclass[10pt,a4paper]{amsart}
\usepackage[margin=19mm]{geometry}
\usepackage{amsmath,amssymb,mathtools}
\usepackage[scr=rsfs,cal=euler]{mathalfa}
\usepackage{xfrac}
\usepackage[unicode,hidelinks,bookmarks=false]{hyperref}
\newcommand{\ie}{i.e.\ }
\newcommand{\cf}{cf.\ }
\newcommand{\resp}{respectively\ }
\newcommand{\Subsection}{\subsection}
\providecommand{\nobreakdash}{\nobreak\hbox{-}}
\setlength{\emergencystretch}{2em}
\multlinegap0pt
\usepackage{lmodern}
\usepackage{mathtools}
\usepackage{enumitem}
\allowdisplaybreaks
\numberwithin{equation}{section}
\setlength{\parindent}{1em}
\setlength{\parskip}{0pt}
\newtheorem{theorem}{Theorem}[section]
\newtheorem{lemma}[theorem]{Lemma}
\newtheorem{proposition}[theorem]{Proposition}
\newtheorem{corollary}[theorem]{Corollary}
\newtheorem{definition}[theorem]{Definition}
\newtheorem{conjecture}[theorem]{Conjecture}
\newtheorem{observation}[theorem]{Observation}
\newtheorem{remark}{Remark}[section]
\newcommand{\Prob}{\mathbb{P}}
\newcommand{\Ex}{\mathbb{E}}
\newcommand{\Use}{\operatorname{Use}}
\newcommand{\Bin}{\operatorname{Bin}}
\newcommand{\deltaq}{\delta_q}
\newcommand{\ee}{\mathrm{e}}
\newcommand{\equalparts}[2]{\{#1\star #2\}}
\begin{document}
\title[Codegree Thresholds for lambda-Choosability of Graphs]{Codegree Thresholds for $\lambda$-Choosability of Graphs}
\author{Chunqiu Fang, Rongxing Xu}
\date{}
\maketitle
\noindent\emph{Source and licence.} Codegree Thresholds for $\lambda$-Choosability of Graphs. Version arXiv:2608.16740v1 (CC BY 4.0). This material is distributed under the Creative Commons Attribution 4.0 International licence, http://creativecommons.org/licenses/by/4.0/, as recorded in the arXiv OAI licence metadata for this exact version and shown on the abstract page https://arxiv.org/abs/2608.16740v1. Full rights notice: licenses/arxiv-2608.16740-CC-BY-4.0.txt.\par\smallskip
\noindent\textit{Editorial scope.} Counted unit: the original \texttt{theorem} environment \texttt{thm:existence} at source lines 191--197 together with its complete proof at lines 1127--1127; the delivered document keeps source lines 112--1180 so that every referenced label and every citation resolves inside it. Native editable source is the paper's own concatenated TeX; the AMS wrapper is editorial formatting only. No publisher class, upstream script or unrelated artwork is redistributed.
\section{Introduction}

All graphs in this paper are finite and simple. We refer to \cite{BondyMurty2008} for graph-theoretic terminology and notation not defined here. For a positive integer $k$, let $[k]=\{1,2,\ldots,k\}$.

List coloring was introduced independently by Vizing~\cite{Vizing1976} and
Erd\H{o}s, Rubin and Taylor~\cite{ERT1979}. A list assignment $L$ of a graph
$G$ assigns a set $L(v)$ of colors to each vertex $v$. An $L$-coloring is a
proper coloring $\phi$ such that $\phi(v) \in L(v)$ for every vertex $v$. For a
positive integer $k$, the graph $G$ is $k$-choosable if it has an $L$-coloring
for every list assignment $L$ with $|L(v)| = k$ for every vertex $v$. The choice
number $\operatorname{ch}(G)$ is the minimum such $k$.

The chromatic number and the choice number can be very different. The complete
bipartite graph $K_{n,n}$ is 2-colorable, while
$\operatorname{ch}(K_{n,n}) = (1 + o(1))\log_2 n$ as $n \to \infty$
\cite{ERT1979}. Alon~\cite{Alon1992} proved more generally that the choice
number of the complete $r$-partite graph with $n$ vertices in each part has
order $r\log_2 n$. He later proved that every graph with minimum degree $d$ has choice
number at least $(1/2 - o(1))\log_2 d$~\cite{Alon2000}. Saxton and
Thomason~\cite{ST2015} introduced the hypergraph container method and used it
to replace the constant $1/2$ by the best possible constant $1$. The same
method was introduced independently by Balogh, Morris and
Samotij~\cite{BMS2015}. Thus the choice number tends to infinity with the
minimum degree, although the chromatic number need not do so.

Several weaker forms of list coloring have also been studied.
Kratochv{\'\i}l, Tuza and Voigt~\cite{KTV1998Complexity} introduced separation
choosability, in which the lists of adjacent vertices are required to have
only a small intersection. Kr\'al' and Sgall~\cite{KralSgall2005} introduced
list coloring with a bounded palette, in which all lists use a common palette
of bounded size. For Alon-type results on these variants, especially for
bipartite graphs and graphs of large minimum degree, see
\cite{KTV1998Brooks,FKK2014,EKT2019,Kang2013,BonamyKang2017}.

In this paper, we consider another variant of list coloring, called
$\lambda$-choosability, which was introduced recently by
Zhu~\cite{Zhu2020}. Let
$\lambda = \{k_1,\ldots,k_q\}$ be a partition, where the $k_i$ need not be distinct, and let $|\lambda| = k_1 + \cdots + k_q$. A
$|\lambda|$-list assignment $L$ is a $\lambda$-assignment if its color set can
be partitioned into $q$ pairwise disjoint sets $X_1,\ldots,X_q$ such that
$|L(v) \cap X_i| = k_i$ for every vertex $v$ and every $i \in [q]$. The graph $G$ is
$\lambda$-choosable if it is $L$-colorable for every $\lambda$-assignment
$L$. This notion puts ordinary coloring and list coloring in the same
framework. For positive integers $a$ and $b$, let $\{a \star b\}$ denote the
partition with $b$ parts, all equal to $a$. The cases $\{k\}$ and
$\{1 \star k\}$ are ordinary $k$-choosability and ordinary
$k$-colorability, respectively. This notion has been studied extensively in recent years. See
\cite{GuZhu2023,ZhuZhu2021,GuJiangWoodZhu2023,ZhuZhu2025} for recent work.

It is natural to ask whether an analogue of Alon's result holds for
$\lambda$-choosability. For a fixed partition $\lambda$, does every graph of
sufficiently large minimum degree fail to be $\lambda$-choosable? The answer
is negative when $\lambda$ has more than one part. For example, $K_{n,n}$ has
minimum degree $n$ and is $\{1,1\}$-choosable. More generally, if $\lambda$
has $q \geq 2$ parts, then the complete $q$-partite graph with $n$ vertices in
each part has minimum degree $(q - 1)n$ and is $q$-colorable, and hence
$\lambda$-choosable. Thus a graph may be
$\lambda$-choosable even when its minimum degree is arbitrarily large. We
prove that a graph is not $\lambda$-choosable if every set of
$q$ vertices has sufficiently many common neighbors.

Formally, for a graph $G$ and a set $S \subseteq V(G)$, let
$N_G(S) = \bigcap_{v \in S} N_G(v)$, where
$N_G(\varnothing)=V(G)$. For an integer $t$ with
$1\leq t\leq |V(G)|$, the
\emph{minimum $t$-codegree} of $G$ is defined by
\[
 \delta_t(G)
 = \min\bigl\{|N_G(S)| : S \subseteq V(G),\ |S|=t\bigr\}.
\]
When the graph is clear, we use $N(S)$ for $N_G(S)$. If $|S| \leq t$, then
$|N_G(S)| \geq \delta_t(G)$.

For a partition $\lambda=\{k_1,\ldots,k_q\}$, let $f(\lambda)$ be the least
integer $d$ such that every graph $G$ with $|V(G)|\geq q$ and
$\delta_q(G)\geq d$ is not $\lambda$-choosable, and let
$f(\lambda)=\infty$ if no such integer exists.
Our first main result is the following.


\begin{theorem}\label{thm:existence}
Let $\lambda=\{k_1,\ldots,k_q\}$ be a partition. There exists an integer
$D_\lambda$ such that every graph $G$ with $\delta_q(G)\geq D_\lambda$
is not $\lambda$-choosable, that is, $f(\lambda)<\infty$. Moreover,
$f(\lambda)\leq 2^{(2q+o(1))|\lambda|}$  as $|\lambda|\to\infty$
whenever $q$ is fixed.
\end{theorem}


For positive integers $r$ and $n$, let $T_r(n)$ denote the complete
$r$-partite graph with $n$ vertices in each part. The $q$-codegree in
Theorem~\ref{thm:existence} cannot be replaced by $p$-codegree for any $p < q$.
Indeed, the graph $T_q(n)$ is
$q$-colorable and hence $\lambda$-choosable for every partition $\lambda$ with
$q$ parts, while $\delta_p(T_q(n)) = (q - p)n$ for $p < q$.
For each positive integer $q$, let $\rho_q$ denote the unique $x\in(0,1)$
satisfying $x=(1-x)^q$. This number is well defined since
$(1-x)^q-x$ is strictly decreasing from $1$ to $-1$ on $[0,1]$.
We obtain the following two lower bounds from complete $(q+1)$-partite graphs.

\begin{proposition}\label{prop:lower-threshold}
Let $\lambda=\{k_1,\ldots,k_q\}$ be a partition. If $N_1$ is a positive integer satisfying
$N_1(q+1)(q/(q+1))^{|\lambda|}<1$,
then $T_{q+1}(N_1)$ is $\lambda$-choosable. Consequently,
\[
f(\lambda)\geq\left\lceil(q+1)^{-1}(1+\frac{1}{q})^{|\lambda|}\right\rceil.
\]
Moreover, if $k$, $q$ and $N_2$ are positive integers satisfying
$(q+1)N_2\rho_q^k<1$,
then $T_{q+1}(N_2)$ is $\equalparts{k}{q}$-choosable. Consequently,
\[
f(\equalparts{k}{q})\geq\left\lceil \frac{\rho_q^{-k}}{(q+1)}\right\rceil.
\]
\end{proposition}

By Theorem~\ref{thm:existence} and the first part of
Proposition~\ref{prop:lower-threshold}, for every fixed $q$ we have
\[
 2^{(\log_2(1+1/q)-o(1))|\lambda|}
 \leq f(\lambda)
 \leq 2^{(2q+o(1))|\lambda|}
\]
as
$|\lambda|\to\infty$. By the equal-part case of
Proposition~\ref{prop:lower-threshold}, we obtain a stronger lower bound. We
then establish a matching upper bound.

\begin{theorem}\label{thm:equal-upper}
Let $q\geq 1$ be fixed. We have
$f(\equalparts{k}{q}) \leq\rho_q^{-(1+o(1))k}$ as $k\to\infty$.
\end{theorem}

By Theorem~\ref{thm:equal-upper} and the second part of
Proposition~\ref{prop:lower-threshold}, we obtain the asymptotic threshold.

\begin{corollary}\label{cor:equal-threshold}
For every fixed $q\geq 1$, $f(\equalparts{k}{q})=\rho_q^{-(1+o(1))k}$ as $k\to\infty$. 
\end{corollary}

We also use the equal-part result to study a corresponding choice parameter.
For a graph $G$, let its \emph{$q$-partition choice number} be
\[
 \operatorname{ch}^{(q)}(G)
 =\min\{k:G\text{ is }\equalparts{k}{q}\text{-choosable}\}.
\]
Thus $\operatorname{ch}^{(1)}(G)=\operatorname{ch}(G)$, and
$\operatorname{ch}^{(q)}(G)=1$ if and only if $G$ is $q$-colorable.

\begin{corollary}\label{cor:partition-choice}
For every fixed $q\geq1$,
\[
 \operatorname{ch}^{(q)}(G)\geq
 \left(\frac{1}{\ln(1/\rho_q)}-o(1)\right)\ln\delta_q(G)
\]
as $\delta_q(G)\to\infty$. The coefficient
$1/\ln(1/\rho_q)$ is asymptotically best possible.
\end{corollary}

\begin{proof}
Let $d=\delta_q(G)$ and fix a sufficiently small $\varepsilon>0$. Take
$k=\lfloor(1-\varepsilon)\ln d/\ln(1/\rho_q)\rfloor$.
Then $k\to\infty$ as $d\to\infty$. By Corollary~\ref{cor:equal-threshold}, $f(\equalparts{k}{q})\leq \rho_q^{-(1+\varepsilon/2)k}$ for sufficiently large $d$.
Since $k\ln(1/\rho_q)\leq(1-\varepsilon)\ln d$, we have
\[
 \begin{aligned}
 \ln f(\equalparts{k}{q}) \leq \left(1+\frac{\varepsilon}{2}\right)
 k\ln(1/\rho_q) \leq \left(1+\frac{\varepsilon}{2}\right)
 (1-\varepsilon)\ln d
 <\ln d.
 \end{aligned}
\] 
Thus $f(\equalparts{k}{q})<d$, so $G$ is not
$\equalparts{k}{q}$-choosable. It follows that
\[
 \operatorname{ch}^{(q)}(G)\geq k+1
 >(1-\varepsilon)\frac{\ln d}{\ln(1/\rho_q)}.
\]
The lower bound follows since $\varepsilon$ can be chosen arbitrarily small.

To see that the coefficient is best possible, let
$N=\lceil\rho_q^{-k}/(q+1)\rceil-1$ and let $G=T_{q+1}(N)$.
By the definition of $N$, we have $(q+1)N\rho_q^k<1$ whenever $N$ is positive,
and $N$ is positive for all sufficiently large $k$.
By Proposition~\ref{prop:lower-threshold}, the graph $G$ is
$\equalparts{k}{q}$-choosable. Hence
$\operatorname{ch}^{(q)}(G)\leq k$, while $\delta_q(G)=N$. Since
$\ln N=(1+o(1))k\ln(1/\rho_q)$,
\[
 \operatorname{ch}^{(q)}(G)
 \leq\left(\frac{1}{\ln(1/\rho_q)}+o(1)\right)
 \ln\delta_q(G).
\]
The coefficient is therefore asymptotically best possible.
\end{proof}

When $q=1$, we have $\rho_1=1/2$,
$\operatorname{ch}^{(1)}(G)=\operatorname{ch}(G)$, and
$\delta_1(G)=\delta(G)$. Corollary~\ref{cor:partition-choice} therefore
becomes $\operatorname{ch}(G)\geq(1-o(1))\log_2\delta(G)$. By the standard
minimum-degree subgraph argument, every graph $G$ of average degree $d$
contains a subgraph $H$ with $\delta(H)\geq d/2$. Hence
$\operatorname{ch}(G)\geq\operatorname{ch}(H)
\geq(1-o(1))\log_2(d/2)=(1-o(1))\log_2d$. This is the asymptotically sharp
lower bound in terms of average degree proved by Saxton and
Thomason~\cite{ST2015}.

The rest of the paper is organized as follows. In
Section~\ref{sec:complete-multipartite}, we prove
Proposition~\ref{prop:lower-threshold}. In Section~\ref{sec:surrounded}, we
prove Theorem~\ref{thm:existence}.
In Section~\ref{sec:partition}, we prove Theorem~\ref{thm:equal-upper}.
 
\section{Proof of Proposition~\ref{prop:lower-threshold}}
\label{sec:complete-multipartite}
 
For every positive integer $N$, we have $\delta_q(T_{q+1}(N))=N$.
In each of the two cases below, let $V_0,V_1,\ldots,V_q$ denote the parts of
the complete $(q+1)$-partite graph under consideration.

We use the following observation in both cases. Suppose that the colors are
labeled by $0,1,\ldots,q$ and that every vertex in $V_j$ has a color in its
list labeled $j$. Color each vertex with one such color. A color has only one
label, so vertices in different parts cannot receive the same color. Thus the
coloring is proper.

We first consider the general partition $\lambda$. Let $L$ be an arbitrary
$\lambda$-assignment of $T_{q+1}(N_1)$, and label every color independently
and uniformly from $\{0,1,\ldots,q\}$. For a vertex $v\in V_j$, the
probability that $L(v)$ contains no color labeled $j$ is
$(q/(q+1))^{|\lambda|}$.
Hence the expected number of vertices having no color labeled with the index
of their part is
\[
 (q+1)N_1\left(\frac{q}{q+1}\right)^{|\lambda|}<1.
\]
Some labeling leaves no such vertex. By the preceding observation, we obtain
a proper $L$-coloring. Thus $T_{q+1}(N_1)$ is $\lambda$-choosable.

To obtain the stated lower bound for $f(\lambda)$, take
\[
 N_1=
 \left\lceil
 \frac{1}{q+1}\left(1+\frac1q\right)^{|\lambda|}
 \right\rceil-1
\]
whenever this integer is positive. This value satisfies the strict inequality
in the proposition, and the graph above has $q$-codegree $N_1$. Hence
\[
 f(\lambda)\geq N_1+1
 =\left\lceil
 \frac{1}{q+1}\left(1+\frac1q\right)^{|\lambda|}
 \right\rceil.
\]
If $N_1=0$, the right-hand side is $1$. In this case the same conclusion
follows from the $\lambda$-choosable graph $K_q$, for which
$\delta_q(K_q)=0$.

We now assume that every part of the partition is equal to $k$. Let $L$ be an
arbitrary $\equalparts{k}{q}$-assignment of $T_{q+1}(N_2)$, with color pools
$X_1,\ldots,X_q$. For every $i\in[q]$, independently assign each color in
$X_i$ label $0$ with probability $\rho_q$ and label $i$ with probability
$1-\rho_q$. If $v\in V_0$, the probability that $L(v)$ contains no color
labeled $0$ is $(1-\rho_q)^{qk}=\rho_q^k$.
If $v\in V_i$ for some $i\in[q]$, only the $k$ colors of $L_i(v)$  
may receive label $i$, and the probability that none of them does is
$\rho_q^k$.
The expected number of vertices having no color labeled with the index of
their part is therefore $(q+1)N_2\rho_q^k<1$.
Some labeling again leaves no such vertex. By the preceding observation,
we obtain a proper $L$-coloring, so $T_{q+1}(N_2)$ is
$\equalparts{k}{q}$-choosable.

Finally, take $N_2=\left\lceil \rho_q^{-k}/(q+1)\right\rceil-1$
whenever this integer is positive. Then $(q+1)N_2\rho_q^k<1$, and therefore
\[
 f(\equalparts{k}{q})\geq N_2+1
 =\left\lceil\frac{\rho_q^{-k}}{q+1}\right\rceil.
\]
If $N_2=0$, the right-hand side is again $1$, and the conclusion follows from
$K_q$.


\section{Proof of Theorem~\ref{thm:existence}}
\label{sec:surrounded}

Let $\lambda=\{k_1,k_2,\ldots,k_q\}$ be a fixed partition. Let $G$ be a
graph, and let $d=\delta_q(G)$.
Let $X_1,\ldots,X_q$ be pairwise disjoint color pools with $|X_i| = k_i^2$ for every $i \in [q]$.
 

The proof uses a deferred-exposure form of Alon's random-list
argument~\cite{Alon2000}. The ordinary list-coloring argument works as
follows. All lists use a fixed color palette and are chosen independently at
random at the beginning, but are revealed only when needed. The goal is to
find a list assignment $L$ for which $G$ has no $L$-coloring. It is enough
to show that this occurs with positive probability. We first reveal
the lists on a sparse set $B$. Many vertices outside $B$ are then
\emph{good}, each of which has the property that for every $L$-coloring of
$G[B]$, the colors used by its $B$-neighbors meet every half-set of the
palette. These colors therefore contain more than half of the palette. When
the previously unexposed random list at a good vertex is finally revealed,
it has a substantial chance of containing only colors already used by its
$B$-neighbors. We complete the argument with a union bound over the
$L$-colorings of $G[B]$.

For a general partition $\lambda=\{k_1,\ldots,k_q\}$, the color palette is
divided into $q$ pools. To play the same role as a good vertex in the
ordinary list-coloring argument, a vertex must satisfy the half-set property
in every pool for every $L$-coloring of $G[B]$. It is difficult to establish
all these properties by revealing $B$ in one step. We therefore divide $B$
into layers and reveal them in stages. By revealing the layers in this order,
we establish the required half-set properties recursively. Before describing the random experiment that chooses the lists and assigns
vertices to the layers, we introduce the notation and parameters used in the
construction.



For $i \in [q]$ and $\varnothing \ne I \subseteq [q]$, let
\[
 \alpha_i =
 \frac{\binom{\lfloor |X_i|/2\rfloor + 1}{k_i}}
 {\binom{|X_i|}{k_i}}=\frac{\binom{\lfloor k_i^2/2\rfloor + 1}{k_i}}
 {\binom{k_i^2}{k_i}},
 \quad
 \rho_I =
 \prod_{i \in I}
 \frac{\binom{\lceil |X_i|/2\rceil}{k_i}}
 {\binom{|X_i|}{k_i}}=\prod_{i \in I}
 \frac{\binom{\lceil k_i^2 /2\rceil}{k_i}}
 {\binom{k_i^2}{k_i}}.
\]
Let $c_{[q]}=1$. For $\varnothing\ne I\subsetneq[q]$, define $c_I$ and
$\hat{c}_I$ recursively by
\begin{equation}\label{eq:c-recursion}
 \hat{c}_{I} = \sum_{J \supsetneq I}c_J,
 \quad
 c_I = \frac{12\hat{c}_{I}\ln |\lambda|}
 {\rho_I\prod_{j \notin I}\alpha_j}.
\end{equation}
Finally, let
$\widetilde{c} = \sum_{\varnothing \ne J \subseteq [q]}c_J$.

With these parameters fixed, assume for the moment that
$\widetilde{c}/\sqrt{d} \leq 1/8$.
The random experiment consists of the following two independent families of
choices.
\begin{enumerate}
\item For every vertex $v$ and every $i \in [q]$, choose $L_i(v)$ uniformly
from $\binom{X_i}{k_i}$, independently over all pairs $(v,i)$, and define
$L(v)=\bigcup_{i=1}^q L_i(v)$.
\item Independently of all list components, assign each vertex $v$ a label
$I\subseteq[q]$, independently over the vertices. Let $B_I$ be the set of
vertices with label $I$. For every nonempty $I\subseteq[q]$, let
$\Prob(v\in B_I)=c_I/\sqrt d$, and let
$\Prob(v\in B_\varnothing)=1-\widetilde c/\sqrt d$.
\end{enumerate}
For every $I \subseteq [q]$, define
\[
 \hat{B}_{I} = \bigcup_{J \supsetneq I}B_J.
\]
Thus $\hat{B}_{[q]}=\varnothing$, while $\hat{B}_{\varnothing}$ is the union
of all nonempty layers. Let $B=\hat{B}_{\varnothing}$, as in Alon's
argument~\cite{Alon2000}. Since the sets $B_I$ form a partition of $V(G)$,
we have $B_\varnothing=V(G)\setminus B$.

\subsection{Blocking events and surrounded vertices}\label{subsec:blocking-events}

A \emph{state} is a triple $(S,I,\mathbf{P}_I)$ satisfying
$S \subseteq V(G)$,
$I \subseteq [q]$, $|S| = |I|$, and
$\mathbf{P}_I = (P_i)_{i \in I}$, where
$P_i\in\binom{X_i}{\lceil |X_i|/2\rceil}$ for every $i\in I$.
The empty state has $S = I = \varnothing$ and the unique empty tuple
$\mathbf{P}_\varnothing$.
For $I \subseteq [q]$, a proper $L$-coloring $\phi$ of
$G[\hat{B}_{I} \cap N(S)]$, a vertex $u \in N(S)$, and $j \notin I$, define
\[
 \Use_j^\phi(u, S, I) =
 \phi\bigl(\hat{B}_{I} \cap N(S) \cap N(u)\bigr) \cap X_j.
\]
We say that $\phi$ \emph{avoids} $\mathbf{P}_I$ if it uses no color in
$\bigcup_{i \in I}P_i$.

We first define the obstruction needed at a fixed state.

\begin{definition}\label{def:blocking-events}
For a state $(S,I,\mathbf{P}_I)$ with $I \ne \varnothing$, let
$\mathcal{E}(S,I,\mathbf{P}_I)$ be the event that, for every proper
$L$-coloring $\phi$ of
$G[\hat{B}_{I} \cap N(S)]$ that avoids $\mathbf{P}_I$, there is a vertex
$u \in B_I \cap N(S)$ satisfying
\begin{enumerate}[label=(\alph*),ref=(\alph*)]
\item\label{item:event-prescribed}
$L_i(u) \subseteq P_i$ for every $i \in I$,
\item\label{item:event-blocking}
$L_j(u) \subseteq \Use_j^\phi(u, S, I)$ for every $j \notin I$.
\end{enumerate}
We call the state $(S,I,\mathbf{P}_I)$ \emph{blocking} if
$\mathcal{E}(S,I,\mathbf{P}_I)$ occurs.
\end{definition}

When $I = [q]$, Definition~\ref{def:blocking-events} has the following
particularly simple form. Since $\hat{B}_{[q]} = \varnothing$, the graph
$G[\hat{B}_{[q]} \cap N(S)]$ has the unique empty coloring $\phi$, which
avoids $\mathbf{P}_{[q]}$. Moreover, there is no $j \notin [q]$, so
Condition~\ref{item:event-blocking} is vacuous. Hence
$\mathcal{E}(S,[q],\mathbf{P}_{[q]})$ is precisely the event that there is a
vertex $w \in B_{[q]} \cap N(S)$ such that
$L_i(w) \subseteq P_i$ for every $i \in [q]$.

The role of the blocking event is to prevent a proper $L$-coloring $\phi$
of $G[\hat{B}_I\cap N(S)]$ that avoids $\mathbf{P}_I$ from being extended
to the corresponding vertex $u$ while still avoiding $\mathbf{P}_I$.
Condition~\ref{item:event-blocking} prevents $u$ from using any list
component $L_j(u)$ with $j\notin I$, since every color in $L_j(u)$ is
already used by a neighbor of $u$ under $\phi$. Thus every proper extension
of $\phi$ to $u$ must use a color from $L_i(u)$ for some $i\in I$. By
Condition~\ref{item:event-prescribed}, $L_i(u)\subseteq P_i$. The extension
therefore uses a color in $\bigcup_{i\in I}P_i$ and no longer avoids
$\mathbf{P}_I$.
 
To find the vertex $u$ required in
Definition~\ref{def:blocking-events}, we need the following definition,
where a member of
$\binom{X_i}{\lceil |X_i|/2\rceil}$ is called a \emph{half-set} of $X_i$
for every $i\in[q]$.

\begin{definition}\label{def:surrounded}
Let $(S,I,\mathbf{P}_I)$ be a state with $I \subsetneq [q]$. A vertex
$u \in N(S) \setminus \hat{B}_{I}$ is \emph{surrounded at}
$(S,I,\mathbf{P}_I)$ if, for every $j \notin I$ and every half-set $P_j$ of
$X_j$, the blocking event
$\mathcal{E}(S \cup \{u\},J,\mathbf{P}_J)$ occurs, where
$J = I \cup \{j\}$ and $\mathbf{P}_J$ is the tuple extending $\mathbf{P}_I$ with
$j$-coordinate $P_j$.
\end{definition}

The role of surrounded vertices is captured by the following lemma. In every
remaining color pool, the neighbors of a surrounded vertex use more than half
of the colors.

\begin{lemma}\label{lem:recursive-hit}
Suppose that $u$ is surrounded at
$(S,I,\mathbf{P}_I)$, where $I$ is allowed to be empty, and let $\phi$ be a
proper $L$-coloring of $G[\hat{B}_{I} \cap N(S)]$ that avoids
$\mathbf{P}_I$.
For every $j \notin I$, the set $\Use_j^\phi(u, S, I)$ meets every
half-set of $X_j$. Consequently,
\[
 |\Use_j^\phi(u, S, I)| \geq \lfloor k_j^2/2\rfloor + 1.
\]
\end{lemma}

\begin{proof}
Fix $j \notin I$ and suppose that some half-set $P_j$ of $X_j$ is disjoint
from $\Use_j^\phi(u, S, I)$. Let $J = I \cup \{j\}$, and let $\mathbf{P}_J$
extend $\mathbf{P}_I$ by the $j$-coordinate $P_j$.
Since $u$ is surrounded at $(S,I,\mathbf{P}_I)$, the blocking event
$\mathcal{E}(S \cup \{u\},J,\mathbf{P}_J)$ occurs.

Let $\phi'$ be the restriction
of $\phi$ to $\hat{B}_{J} \cap N(S \cup \{u\})$. Since $J \supsetneq I$, we have
$\hat{B}_{J} \cap N(S \cup \{u\})
 \subseteq \hat{B}_{I} \cap N(S) \cap N(u)$. Thus $\phi'$ is a proper
$L$-coloring of $G[\hat{B}_{J} \cap N(S \cup \{u\})]$ and avoids
$\mathbf{P}_I$. We claim that $\phi'$ also uses no color in $P_j$. If
$\phi'(v) \in P_j$ for some
$v \in \hat{B}_{J} \cap N(S \cup \{u\})$, then
$v \in \hat{B}_{I} \cap N(S) \cap N(u)$ and $\phi'(v) = \phi(v)$. Since
$P_j \subseteq X_j$, by the definition of $\Use_j^\phi(u, S, I)$, we have
$\phi'(v) \in \Use_j^\phi(u, S, I)$. Hence
$\phi'(v) \in P_j \cap \Use_j^\phi(u, S, I)$, contrary to the choice of $P_j$.
Therefore, $\phi'$ avoids $\mathbf{P}_J$.

Since the blocking event $\mathcal{E}(S \cup \{u\},J,\mathbf{P}_J)$ occurs and
$\phi'$ avoids $\mathbf{P}_J$, there is a vertex
$z \in B_J \cap N(S \cup \{u\})$ by
Definition~\ref{def:blocking-events}. In particular,
$z \in \hat{B}_{I} \cap N(S) \cap N(u)$, so $\phi(z)$ is defined. Choose
$r \in [q]$ such that $\phi(z) \in L_r(z)$. If $r \in I$, then by
Definition~\ref{def:blocking-events}\ref{item:event-prescribed}, we have
$\phi(z) \in P_r$, contrary to the fact that $\phi$ avoids $\mathbf{P}_I$.
If $r = j$, then by
Definition~\ref{def:blocking-events}\ref{item:event-prescribed} and the definition of
$\Use_j^\phi(u, S, I)$, we have
$\phi(z) \in P_j \cap \Use_j^\phi(u, S, I)$, contrary to the choice of $P_j$.
Finally, if $r \notin J$, then by
Definition~\ref{def:blocking-events}\ref{item:event-blocking}, we have
$\phi(z) \in \Use_r^{\phi'}(z, S \cup \{u\}, J)$. Hence there is a vertex
$v \in \hat{B}_{J} \cap N(S \cup \{u\}) \cap N(z)$ such that
$\phi(v) = \phi'(v) = \phi(z)$, contradicting the fact that $\phi$ is a proper coloring on $G[\hat{B}_{I} \cap N(S)]$.

Thus $\Use_j^\phi(u, S, I)$ meets every
half-set of $X_j$. Its complement has size at most
$\lceil k_j^2/2\rceil - 1$, and hence
$|\Use_j^\phi(u, S, I)| \geq \lfloor k_j^2/2\rfloor + 1$.
\end{proof}
 

\subsection{The local blocking lemma}\label{subsec:local-blocking}

The blocking events were defined in the previous subsection. We now prove that,
for every state $(S,I,\mathbf{P}_I)$ with $I \ne \varnothing$, the probability
that the blocking event $\mathcal{E}(S,I,\mathbf{P}_I)$ does not occur is at most
$\varepsilon_I(d)$, where the error functions $\varepsilon_I(d)$ are defined
recursively as follows. First, let
\[
 \varepsilon_{[q]}(d) = \exp(-\rho_{[q]}\sqrt{d}).
\]
For $\varnothing \ne I \subsetneq [q]$, once $\varepsilon_J(d)$ has been
defined for every $J \supsetneq I$, let
\begin{align*}
 \varepsilon_I(d)
=  \exp\left(-\frac{\hat{c}_{I}\sqrt{d}}{3}\right)
 +4\sum_{j \notin I}
 \binom{k_j^2}{\lceil k_j^2/2\rceil}
 \varepsilon_{I \cup \{j\}}(d) 
 +\exp\left(-\frac{c_I\rho_I\sqrt{d}}{16}\right)
 +\exp\left(-\hat{c}_{I}\sqrt{d}\ln |\lambda|\right).
\end{align*}
For every nonempty $I \subseteq [q]$, by induction on $q - |I|$, we have
$\varepsilon_I(d) \to 0$ as $d \to \infty$.

Before stating and proving the lemma, we introduce the following standard form of the Chernoff bound.

\begin{lemma}\label{lem:chernoff}
	Let $X \sim \Bin(n,p)$ and $\mu=np$. Then 
	\begin{enumerate}[label=(\arabic*),ref=(\arabic*)]
		\item\label{item:chernoff-lower}
		$\Prob(X < \mu/2) \leq e^{-\mu/8}$.
		\item\label{item:chernoff-upper}
		$\Prob(X > 2\mu) \leq e^{-\mu/3}$.
	\end{enumerate}
\end{lemma}

The desired probability bound is as follows.

\begin{lemma}
\label{lem:local-blocking}
Suppose that $d = \delta_q(G)$ satisfies $\widetilde{c}/\sqrt{d} \leq 1/8$.
For every state $(S,I,\mathbf{P}_I)$ with $\varnothing \ne I \subseteq [q]$, the probability
that the blocking event $\mathcal{E}(S,I,\mathbf{P}_I)$ does not occur is at most
$\varepsilon_I(d)$.
\end{lemma}

\begin{proof}
We use induction on $q - |I|$. By the definition of $d = \delta_q(G)$, we have
$|N(S)| \geq d$.

If $I = [q]$, then a fixed vertex $w \in N(S)$ belongs to $B_{[q]}$ and
satisfies $L_i(w) \subseteq P_i$ for every $i \in [q]$ with probability
$\rho_{[q]}/\sqrt d$. These events are independent for distinct vertices.
Hence the probability that the blocking event
$\mathcal{E}(S,[q],\mathbf{P}_{[q]})$ does not occur is at most
\[
 \left(1 - \frac{\rho_{[q]}}{\sqrt d}\right)^{|N(S)|}
 \leq \exp\left(-\frac{\rho_{[q]}|N(S)|}{\sqrt d}\right)
 \leq \exp(-\rho_{[q]}\sqrt{d})
 = \varepsilon_{[q]}(d).
\]

Now let $\varnothing \ne I \subsetneq [q]$ and suppose that the lemma holds
for every $J \supsetneq I$.
We divide the induction step into three steps. In Step~1, we reveal the sets
$B_J$, for $J \supsetneq I$, and all list components of the vertices in
$\hat{B}_{I}$, and show that many vertices
$u \in N(S) \setminus \hat{B}_{I}$ are surrounded at $(S,I,\mathbf{P}_I)$. In
Step~2, we reveal $B_I$ and the components $L_i(u)$ for $u\in B_I$ and
$i\in I$. We show that many of the surrounded vertices have
$L_i(u) \subseteq P_i$ for every $i \in I$. These are the candidate blocking vertices
for the event
$\mathcal{E}(S,I,\mathbf{P}_I)$. In Step~3, we reveal their remaining list
components. For every proper $L$-coloring $\phi$ of
$G[\hat{B}_{I} \cap N(S)]$ that avoids $\mathbf{P}_I$, we show that some
candidate vertex satisfies Definition~\ref{def:blocking-events}%
\ref{item:event-blocking}. Hence the blocking event
$\mathcal{E}(S,I,\mathbf{P}_I)$ occurs.

\medskip
\noindent
\textbf{Step 1.} Reveal the sets $B_J$, $J \supsetneq I$, and the lists on
$\hat{B}_{I}$.\par\nobreak
\medskip

In this step, we reveal the sets $B_J$ for $J \supsetneq I$,
together with all list components of the vertices in their union
$\hat{B}_{I}$. After these choices are revealed, we know whether each blocking
event whose second coordinate strictly contains $I$ occurs. We therefore also
know whether each $u \in N(S) \setminus \hat{B}_{I}$ is surrounded at
$(S,I,\mathbf{P}_I)$, and this decision does not use whether $u \in B_I$ or any
list component of $u$.

Since the sets $B_J$ are pairwise disjoint, each vertex belongs to
$\hat{B}_{I}$ independently with probability $\hat{c}_{I}/\sqrt{d}$.
Hence $|\hat{B}_{I} \cap N(S)|$ has binomial distribution with mean
$\hat{c}_{I}|N(S)|/\sqrt{d}$.
By Lemma~\ref{lem:chernoff}\ref{item:chernoff-upper}, we have
\begin{equation}\label{eq:upper-tail}
 \Prob\left(
 |\hat{B}_{I} \cap N(S)| > \frac{2\hat{c}_{I}|N(S)|}{\sqrt{d}}
 \right)
 \leq \exp\left(-\frac{\hat{c}_{I}|N(S)|}{3\sqrt{d}}\right)
 \leq \exp\left(-\frac{\hat{c}_{I}\sqrt{d}}{3}\right).
\end{equation}

For each $j \notin I$ and each half-set $P_j$ of $X_j$, let
$J = I \cup \{j\}$, and let $\mathbf{P}_J$ be the tuple extending $\mathbf{P}_I$
with $j$-coordinate $P_j$. Let $F$ be the set of vertices
$u \in N(S) \setminus \hat{B}_{I}$ for which the blocking event
$\mathcal{E}(S \cup \{u\},J,\mathbf{P}_J)$ fails for at least one such pair
$(j,P_j)$. Thus $F$ is determined by the choices revealed in Step~1.
For a fixed $u \in N(S)$, by the induction hypothesis and the union bound over
all such pairs $(j, P_j)$, we have
\[
 \Prob(u \in F) \leq
 \sum_{j \notin I}
 \binom{k_j^2}{\lceil k_j^2/2\rceil}
 \varepsilon_{I \cup \{j\}}(d).
\]
By linearity of expectation and Markov's inequality, we have
\begin{equation}\label{eq:failed-vertex-tail}
 \Prob\left(|F| > \frac{|N(S)|}{4}\right)
 \leq 4\sum_{j \notin I}
 \binom{k_j^2}{\lceil k_j^2/2\rceil}
 \varepsilon_{I \cup \{j\}}(d).
\end{equation}
The events that different vertices belong to $F$ need not be independent,
since only linearity of expectation and Markov's inequality are used.

\medskip
\noindent
\textbf{Step 2.} Reveal $B_I$, together with $L_i(u)$ for every $u\in B_I$
and $i\in I$.
\par\nobreak
\medskip

Condition on any possible result of Step~1 for which
\begin{equation}\label{eq:first-step-bounds}
 |\hat{B}_{I} \cap N(S)|
 \leq \frac{2\hat{c}_{I}|N(S)|}{\sqrt{d}},
 \quad
 |F| \leq \frac{|N(S)|}{4}.
\end{equation}
Let $W=N(S)\setminus(\hat{B}_I\cup F)$. By
Definition~\ref{def:surrounded}, this is precisely the set of vertices in
$N(S)\setminus\hat{B}_I$ that are surrounded at
$(S,I,\mathbf{P}_I)$. Since $\hat{c}_I\leq\widetilde{c}$ and
$\widetilde{c}/\sqrt{d}\leq 1/8$, we have
\[
 |W| \geq |N(S)| - \frac{2\hat{c}_{I}|N(S)|}{\sqrt{d}}
 - \frac{|N(S)|}{4} \geq \frac{|N(S)|}{2}.
\]
For $u \in W$, let $\mathcal{A}_u$ be the event that $u \in B_I$ and
$L_i(u) \subseteq P_i$ for every $i \in I$.
For every $u \in W$, the information revealed in Step~1 includes
$u \notin \hat{B}_{I}$. Moreover, since the sets $B_J$ are pairwise
disjoint, $\mathcal{A}_u$ implies $u \notin \hat{B}_{I}$. In the original
random experiment,
\[
 \Prob(\mathcal{A}_u) = \frac{c_I\rho_I}{\sqrt{d}},
 \quad
 \Prob(u \notin \hat{B}_{I})
 = 1 - \frac{\hat{c}_{I}}{\sqrt{d}}.
\]
It follows that
\[
\begin{aligned}
 \Prob\bigl(\mathcal{A}_u \mid u \notin \hat{B}_{I}\bigr) =
 \frac{
 \Prob\bigl(\mathcal{A}_u
 \cap \{u \notin \hat{B}_{I}\}\bigr)}
 {\Prob(u \notin \hat{B}_{I})} =
 \frac{\Prob(\mathcal{A}_u)}
 {\Prob(u \notin \hat{B}_{I})} =
 \frac{c_I\rho_I}{\sqrt{d} - \hat{c}_{I}}
 \geq \frac{c_I\rho_I}{\sqrt{d}}.
\end{aligned}
\]
For $u\in W$, Step~1 reveals that $u\notin\hat{B}_I$, but it does not
reveal whether $u\in B_I$ or any list component of $u$. Conditional on this
result of Step~1, the events $\mathcal A_u$, $u\in W$, remain independent,
and by the preceding
calculation each occurs with probability
$c_I\rho_I/(\sqrt d-\hat c_I)$.

Let $R$ be the set of vertices $u\in W$ for which $\mathcal A_u$ occurs.
Then $|R|$ has a binomial distribution with parameters $|W|$ and
$c_I\rho_I/(\sqrt d-\hat c_I)$. Thus
\[
\Ex(|R|)
=\frac{c_I\rho_I|W|}{\sqrt d-\hat c_I}
\geq\frac{c_I\rho_I|N(S)|}{2\sqrt d}.
\]
By Lemma~\ref{lem:chernoff}\ref{item:chernoff-lower},
\begin{equation}\label{eq:candidate-tail}
 \Prob\left(
 |R| < \frac{c_I\rho_I|N(S)|}{4\sqrt{d}}
 \right)
 \leq \exp\left(-\frac{c_I\rho_I|N(S)|}{16\sqrt{d}}\right)
 \leq \exp\left(-\frac{c_I\rho_I\sqrt{d}}{16}\right).
\end{equation}

\medskip
\noindent
\textbf{Step 3.} Reveal the remaining list components of the vertices in $R$.
\par\nobreak
\medskip

Condition further on the event that
\begin{equation}\label{eq:candidate-count}
 |R| \geq \frac{c_I\rho_I|N(S)|}{4\sqrt{d}}.
\end{equation}
At this point, the set $B_I$, the components $L_i(u)$ for $u\in B_I$ and
$i\in I$, and the set $R$ have been revealed.
Every vertex of $R$ lies in $B_I \cap N(S)$ and satisfies
Definition~\ref{def:blocking-events}\ref{item:event-prescribed}. It is also
surrounded at $(S,I,\mathbf{P}_I)$ by the definition of $W$. Reveal
$L_j(u)$ for every $u \in R$ and $j \notin I$.

After Step~1, the lists on $\hat{B}_{I} \cap N(S)$ are fixed. The number of
proper $L$-colorings of
$G[\hat{B}_{I} \cap N(S)]$ that avoid $\mathbf{P}_I$ is at most
\begin{equation}\label{eq:number-upper-colorings}
 |\lambda|^{|\hat{B}_{I} \cap N(S)|}
 \leq |\lambda|^{2\hat{c}_{I}|N(S)|/\sqrt{d}}
 = \exp\left(
 \frac{2\hat{c}_{I}|N(S)|\ln |\lambda|}{\sqrt{d}}
 \right).
\end{equation}
Fix such a coloring $\phi$. Lemma~\ref{lem:recursive-hit} shows that
$|\Use_j^\phi(u,S,I)|\geq\lfloor k_j^2/2\rfloor+1$ for every $u\in R$ and
$j\notin I$. By
the definition of $\alpha_j$, a uniformly random member of
$\binom{X_j}{k_j}$ is contained in $\Use_j^\phi(u, S, I)$ with probability
at least $\alpha_j$. For the list components revealed in Step~3, a fixed
vertex $u \in R$ therefore satisfies Condition~\ref{item:event-blocking} in
Definition~\ref{def:blocking-events} with probability at least
$\prod_{j \notin I}\alpha_j$. These events are independent over the vertices
of $R$. By Inequality~\eqref{eq:candidate-count}, the probability that no vertex of $R$
satisfies Condition~\ref{item:event-blocking} in
Definition~\ref{def:blocking-events} is at most
\[
 \left(1 - \prod_{j \notin I}\alpha_j\right)^{|R|}
 \leq \exp\left(-|R|\prod_{j \notin I}\alpha_j\right)
 \leq \exp\left(
 -\frac{c_I\rho_I|N(S)|}{4\sqrt{d}}
 \prod_{j \notin I}\alpha_j
 \right).
\]

Under the conditioning in \eqref{eq:first-step-bounds}
and~\eqref{eq:candidate-count}, we obtain from
Equality~\eqref{eq:c-recursion}, Inequality~\eqref{eq:number-upper-colorings},
and the union bound that the probability that some such coloring has no
vertex in $R$ satisfying
Definition~\ref{def:blocking-events}\ref{item:event-blocking} is at most
\begin{equation}\label{eq:third-step-failure}
\begin{aligned}
|\lambda|^{|\hat{B}_{I} \cap N(S)|}
\left(1 - \prod_{j \notin I}\alpha_j\right)^{|R|}
 & \leq \exp\left(
 \frac{2\hat{c}_{I}|N(S)|\ln |\lambda|}{\sqrt{d}}
 -\frac{c_I\rho_I|N(S)|}{4\sqrt{d}}
 \prod_{j \notin I}\alpha_j
 \right) \\
 & = \exp\left(
 -\frac{\hat{c}_{I}|N(S)|\ln |\lambda|}{\sqrt{d}}
 \right)\\
 & \leq
 \exp\left(-\hat{c}_{I}\sqrt{d}\ln |\lambda|\right).
\end{aligned}
\end{equation}

By construction, $R \subseteq B_I \cap N(S)$, and every vertex of $R$
satisfies Definition~\ref{def:blocking-events}\ref{item:event-prescribed}.
If Inequalities~\eqref{eq:first-step-bounds} and~\eqref{eq:candidate-count}
hold and the event estimated in Inequality~\eqref{eq:third-step-failure} does
not occur, then every coloring in Definition~\ref{def:blocking-events} has a
vertex in $R$ satisfying
Definition~\ref{def:blocking-events}\ref{item:event-blocking}. Thus the
blocking event $\mathcal{E}(S,I,\mathbf{P}_I)$ occurs. Since the estimates in
Steps~2 and~3 apply after conditioning on any earlier information satisfying
the required bounds,
by Inequalities~\eqref{eq:upper-tail}, \eqref{eq:failed-vertex-tail},
\eqref{eq:candidate-tail}, and~\eqref{eq:third-step-failure}, the probability
that this blocking event does not occur is at most
\begin{align*}
 \exp\left(-\frac{\hat{c}_{I}\sqrt{d}}{3}\right)
 +4\sum_{j \notin I}
 \binom{k_j^2}{\lceil k_j^2/2\rceil}
 \varepsilon_{I \cup \{j\}}(d)
 +\exp\left(-\frac{c_I\rho_I\sqrt{d}}{16}\right)
 +\exp\left(-\hat{c}_{I}\sqrt{d}\ln |\lambda|\right)
 = \varepsilon_I(d).
\end{align*}
This completes the induction.
\end{proof}

\subsection{The final exposure}
\label{subsec:final-exposure}

We now use the local estimate to select suitable values for the information
exposed on the nonempty layers. Let $A$ be the set of vertices in
$V(G) \setminus B$ that are
surrounded at the empty state. Thus $a \in A$ precisely when the
blocking event $\mathcal{E}(\{a\},\{i\},\mathbf{P}_{\{i\}})$ occurs for every
$i \in [q]$ and every half-set $P_i$ of $X_i$, where
$\mathbf{P}_{\{i\}} = (P_i)$. By Lemma~\ref{lem:local-blocking} and the union
bound, we have
\begin{equation}\label{eq:bad-a}
 \Prob(a \notin A)
 \leq \frac{\widetilde{c}}{\sqrt{d}}
 +\sum_{i = 1}^q
 \binom{k_i^2}{\lceil k_i^2/2\rceil}
 \varepsilon_{\{i\}}(d).
\end{equation}

Choose $D_\lambda$ large enough so that for every $d \geq D_\lambda$,
\begin{enumerate}[label=(D\arabic*),ref=(D\arabic*),labelindent=2em,
	leftmargin=*, itemsep=0.5em,topsep=0.7em]
\item\label{cond:final-error} $\widetilde{c}/\sqrt{d}
+\sum_i\binom{k_i^2}{\lceil k_i^2/2\rceil}
\varepsilon_{\{i\}}(d) \leq 1/16$.
\item\label{cond:final-union} $16\widetilde{c}\ln |\lambda|/\sqrt{d}
 \leq \prod_{i = 1}^q \alpha_i$.
\end{enumerate}
By condition~\ref{cond:final-error}, we have
$\widetilde{c}/\sqrt{d} \leq 1/8$. Such a choice exists
because all constants depend only on $\lambda$ and all error functions tend
to zero.

Let $d = \delta_q(G) \geq D_\lambda$ and let $n = |V(G)|$. By
Inequality~\eqref{eq:bad-a}, the expected number of vertices outside $A$ is at
most $n/16$, while $\Ex(|B|) = \widetilde{c}n/\sqrt{d}$.
By Markov's inequality, we have
\[
 \Prob(|A| < n/2) \leq \frac{1}{8},
 \quad
 \Prob\left(|B| > \frac{4\widetilde{c}n}{\sqrt{d}}\right) \leq \frac{1}{4}.
\]
Consequently, with positive probability the information exposed so far satisfies
\begin{equation}\label{eq:AB-size}
 |A| \geq n/2,
 \quad
 |B| \leq \frac{4\widetilde{c}n}{\sqrt{d}}.
\end{equation}
Condition on any realization of the exposed information satisfying
\eqref{eq:AB-size}. The lists on
$V(G)\setminus B$ remain unexposed.

Let $\Phi_B$ be the set of proper $L$-colorings of $G[B]$.
If $\Phi_B = \varnothing$, then the full random assignment is already
uncolorable, regardless of the unexposed lists. We may therefore assume that
$\Phi_B \ne \varnothing$.

We now reveal the lists on $A$. Under this conditioning,
the components $L_i(a)$, for $a\in A$ and $i\in[q]$, remain mutually
independent and uniformly distributed in $\binom{X_i}{k_i}$. Indeed, whether
$a$ belongs to $A$ was determined without exposing any component of $L(a)$.

Fix $\phi \in \Phi_B$ and $a \in A$. For every $i \in [q]$, since $a$ is
surrounded at the empty state, by Lemma~\ref{lem:recursive-hit} with
$S = I = \varnothing$, we have
\[
 |\Use_i^\phi(a, \varnothing, \varnothing)|
 \geq \lfloor |X_i|/2\rfloor + 1 = \lfloor k_i^2/2\rfloor + 1.
\]
Hence, by the definition of $\alpha_i$ and independence over $i \in [q]$,
the probability that
$L_i(a) \subseteq \Use_i^\phi(a, \varnothing, \varnothing)$ for every
$i \in [q]$ is at least $\prod_{i = 1}^q \alpha_i$.
On this event, no color in $L(a)$ is available for extending $\phi$ to $a$.
By independence over $a \in A$, the probability that $\phi$ extends to every
vertex of $A$ is at most
\[
 \left(1 - \prod_{i = 1}^q \alpha_i\right)^{|A|}
 \leq \exp\left(-|A|\prod_{i = 1}^q \alpha_i\right).
\]
Moreover, $|\Phi_B| \leq |\lambda|^{|B|}$. By the inequalities
in~\eqref{eq:AB-size}, condition~\ref{cond:final-union}, and the union bound over all
$\phi \in \Phi_B$, the probability that some $\phi$ extends to every vertex of
$A$ is at most
\begin{align*}
 \exp\bigl(|B|\ln |\lambda|
           -|A|\prod_{i = 1}^q \alpha_i\bigr) \leq
 \exp\left(
 \frac{4\widetilde{c}n\ln |\lambda|}{\sqrt{d}}
 -\frac{n}{2}\prod_{i = 1}^q \alpha_i\right)
  \leq \exp\left(-\frac{n}{4}\prod_{i = 1}^q \alpha_i\right) < 1.
\end{align*}
We can therefore choose the previously unexposed lists on $A$ so that no
coloring of $G[B]$ extends to all of $A$. The lists already sampled on
$V(G) \setminus (A \cup B)$ are irrelevant. Hence some outcome of the original
single random experiment is an uncolorable $\lambda$-assignment. Since $G$
was arbitrary with $\delta_q(G)\geq D_\lambda$, it follows that
$f(\lambda)\leq D_\lambda<\infty$. We have proved the first assertion of
Theorem~\ref{thm:existence}.

\subsection{The asymptotic upper bound}
\label{subsec:threshold}

It remains to prove the asymptotic upper bound in
Theorem~\ref{thm:existence}. Throughout this subsection, $q$ is fixed, and
every $o(1)$ term tends to zero as $|\lambda|\to\infty$. We first estimate
the probabilities $\alpha_i$ and $\rho_I$, the recursive
constants $c_I$, and the error terms $\varepsilon_I(d)$.

 
The following elementary estimate is the binomial-ratio bound used in
Alon's proof~\cite{Alon2000}. We include its short proof for completeness.

\begin{lemma}\label{lem:half-binomial-ratio}
	For every positive integer $k$ and every integer
	$s\geq\lceil k^2/2\rceil$, $\binom{s}{k}/\binom{k^2}{k} \geq 2^{-(k+1)}$.
\end{lemma} 
\begin{proof}
	The result is immediate for $k=1$, so assume that $k\geq2$. Since $s\geq k^2/2$, we have
	\[ 
	\frac{\binom{s}{k}}{\binom{k^2}{k}} =  \prod_{i=0}^{k-1}\frac{s-i}{k^2-i} \geq \frac{1}{2^k}\prod_{i=0}^{k-1}\frac{k^2-2i}{k^2-i} =
	\frac{1}{2^k}\prod_{i=0}^{k-1}\left(1-\frac{i}{k^2-i}\right) \geq  \frac{1}{2^k}\left(1-\sum_{i=0}^{k-1}\frac{i}{k^2-k}\right)  =2^{-(k+1)}. 
	\]
	The second inequality uses $\prod_i(1-x_i)\geq1-\sum_i x_i$ for $0\leq x_i\leq1$ and
	$k^2-i\geq k^2-k$.
\end{proof} 

By Lemma~\ref{lem:half-binomial-ratio}, applied to the ratios defining
$\alpha_i$ and $\rho_I$, we have
\begin{equation}\label{eq:alpha-rho-lower}
 \alpha_i\geq2^{-(k_i+1)},
 \quad
 \rho_I\geq2^{-|I|-\sum_{i\in I}k_i},
 \quad
 \rho_I\prod_{j\notin I}\alpha_j\geq2^{-q-|\lambda|}.
\end{equation}

We first estimate $\widetilde c$. For all sufficiently large $|\lambda|$, we
prove by induction on $q-|I|$ that
$c_I\leq(12\cdot 2^{2q+|\lambda|}\ln|\lambda|)^{q-|I|}$ for every nonempty
$I\subseteq[q]$.
Indeed, the case $I=[q]$ follows from $c_{[q]}=1$. Suppose that
$\varnothing\ne I\subsetneq[q]$ and this estimate holds for every $J\supsetneq I$. Since
$q-|J|\leq q-|I|-1$ and there are at most $2^q$ such sets $J$,
\[
 \hat c_I=\sum_{J\supsetneq I}c_J
 \leq2^q
 \bigl(12\cdot 2^{2q+|\lambda|}\ln|\lambda|\bigr)^{q-|I|-1}.
\]
By Equality~\eqref{eq:c-recursion} and Inequality~\eqref{eq:alpha-rho-lower},
$c_I\leq(12\cdot 2^{q+|\lambda|}\ln|\lambda|)\hat c_I$. By combining this
inequality with the preceding estimate, we obtain the required bound for
$c_I$ and complete the induction. Since there are at most
$2^q$ nonempty sets $I\subseteq[q]$,
\begin{equation}\label{eq:ctilde-explicit-bound}
 \widetilde c
 =\sum_{\varnothing\ne I\subseteq[q]}c_I
 \leq2^q
 \bigl(12\cdot 2^{2q+|\lambda|}\ln|\lambda|\bigr)^{q-1}.
\end{equation}


We then estimate $\varepsilon_I(d)$. Recall that the terminal case in its
recursive definition is
$\varepsilon_{[q]}(d)=\exp(-\rho_{[q]}\sqrt d)$. By
Inequality~\eqref{eq:alpha-rho-lower},
\[
 \varepsilon_{[q]}(d)
 \leq\exp(-2^{-q-|\lambda|}\sqrt d)
 \leq\exp\left(-\frac{2^{-q-|\lambda|}\sqrt d}{16}\right).
\]
Now let $\varnothing\ne I\subsetneq[q]$. We have $\hat c_I\geq c_{[q]}=1$ and
$\prod_{j\notin I}\alpha_j\leq1$. By Equality~\eqref{eq:c-recursion},
$c_I\rho_I\geq12\ln|\lambda|\geq1$ for all sufficiently large
$|\lambda|$. Also,
$2^{-q-|\lambda|}\leq1$ and $\ln|\lambda|\geq1$. It follows that
\[
 \max\left\{
 \exp\left(-\frac{\hat c_I\sqrt d}{3}\right),
 \exp\left(-\frac{c_I\rho_I\sqrt d}{16}\right),\quad
 \exp(-\hat c_I\sqrt d\ln|\lambda|)
 \right\}
 \leq\exp\left(-\frac{2^{-q-|\lambda|}\sqrt d}{16}\right).
\]
By the bound
$\binom{k_j^2}{\lceil k_j^2/2\rceil}\leq2^{k_j^2}$ in the remaining term of
that definition, we have
\[
 \varepsilon_I(d)
 \leq3\exp\left(-\frac{2^{-q-|\lambda|}\sqrt d}{16}\right)
 +4\sum_{j\notin I}2^{k_j^2}\varepsilon_{I\cup\{j\}}(d).
\]

The first term on the right is the same for every $I$, and the terminal term
$\varepsilon_{[q]}(d)$ is at most this term. Each substitution adds one new
element to the index set, so after at most $q-|I|\leq q-1$ substitutions,
every resulting term, for some $J\subseteq[q]\setminus I$, is at most
\[
 3\exp\left(-\frac{2^{-q-|\lambda|}\sqrt d}{16}\right)
 \prod_{j\in J}\bigl(4\cdot2^{k_j^2}\bigr).
\]
The same set $J$ may occur more than
once if its elements are added in different orders. For every such $J$,
\[
 \prod_{j\in J}4\cdot2^{k_j^2}
 \leq4^q2^{\sum_jk_j^2}
 \leq2^{|\lambda|^2+2q}.
\]
At each substitution there are at most $q$ choices for the new element, so
the full expansion contains at most $1+q+\cdots+q^q$ terms. Consequently,
\begin{equation}\label{eq:error-asymptotic-bound}
 \varepsilon_I(d)
 \leq2^{|\lambda|^2+O_q(1)}
 \exp\left(-\frac{2^{-q-|\lambda|}\sqrt d}{16}\right)
\end{equation}
for every nonempty $I\subseteq[q]$ and all sufficiently large $|\lambda|$.

\medskip
\noindent

\emph{Proof of the second part of Theorem~\ref{thm:existence}.}

Let
\[
 D_\lambda=\left\lceil
 2^{2q|\lambda|+2|\lambda|^{2/3}}
 \right\rceil.
\]
We show that, for all sufficiently large $|\lambda|$,
conditions~\ref{cond:final-error} and~\ref{cond:final-union}
hold whenever $d\geq D_\lambda$.
If $d\geq D_\lambda$, then
$\sqrt d\geq2^{q|\lambda|+|\lambda|^{2/3}}$.
It follows from Inequality~\eqref{eq:ctilde-explicit-bound} that
\[
 \frac{\widetilde c}{\sqrt d}
 \leq
 \frac{2^q(12\cdot 2^{2q}\ln|\lambda|)^{q-1}}
 {2^{|\lambda|+|\lambda|^{2/3}}}
 =o(1)
\]
and
\[
 \frac{16\widetilde c\ln|\lambda|}{\sqrt d}
 \leq
 \frac{2^{q+4}(12\cdot 2^{2q})^{q-1}(\ln|\lambda|)^q}
 {2^{|\lambda|+|\lambda|^{2/3}}}
 \leq2^{-q-|\lambda|}
 \leq\prod_{i=1}^q\alpha_i
\]
for all sufficiently large $|\lambda|$. We have therefore verified
condition~\ref{cond:final-union}.

The number of half-sets of $X_i$ is at most $2^{k_i^2}$. Hence, by
Inequality~\eqref{eq:error-asymptotic-bound},
\[ 
 \sum_{i=1}^q
 \binom{k_i^2}{\lceil k_i^2/2\rceil}\varepsilon_{\{i\}}(d) \leq2^{2|\lambda|^2+O_q(1)}
 \exp\left(
 -\frac{2^{-q}2^{(q-1)|\lambda|+|\lambda|^{2/3}}}{16}
 \right) =o(1). 
\]
By combining this estimate with $\widetilde c/\sqrt d=o(1)$, we verify
condition~\ref{cond:final-error}.
Therefore the argument in the preceding subsection applies to every graph
$G$ with $\delta_q(G)\geq D_\lambda$. Thus
\[
 f(\lambda)\leq D_\lambda
 =2^{(2q+o(1))|\lambda|}.
\]
This completes the proof of Theorem~\ref{thm:existence}.

\section{\texorpdfstring{\thinspace}{}Proof of
Theorem~\ref{thm:equal-upper}}\label{sec:partition}


\begin{thebibliography}{99}
\bibitem[Alon1992]{Alon1992} N.~Alon, \newblock Choice numbers of graphs: a probabilistic approach, \newblock \emph{Combin. Probab. Comput.} 1 (1992), 107--114.
\bibitem[Alon2000]{Alon2000} N.~Alon, \newblock Degrees and choice numbers, \newblock \emph{Random Structures Algorithms} 16 (2000), 364--368.
\bibitem[BMS2015]{BMS2015} J.~Balogh, R.~Morris and W.~Samotij, \newblock Independent sets in hypergraphs, \newblock \emph{J. Amer. Math. Soc.} 28 (2015), 669--709.
\bibitem[BonamyKang2017]{BonamyKang2017} M.~Bonamy and R.~J. Kang, \newblock List colouring with a bounded palette, \newblock \emph{J. Graph Theory} 84 (2017), 93--103.
\bibitem[BondyMurty2008]{BondyMurty2008} J.~A. Bondy and U.~S.~R. Murty, \newblock \emph{Graph Theory}, \newblock Graduate Texts in Mathematics, vol.~244, Springer, London, 2008.
\bibitem[EKT2019]{EKT2019} L.~Esperet, R.~J. Kang and S.~Thomass\'e, \newblock Separation choosability and dense bipartite induced subgraphs, \newblock \emph{Combin. Probab. Comput.} 28 (2019), 720--732.
\bibitem[ERT1979]{ERT1979} P.~Erd\H{o}s, A.~L. Rubin and H.~Taylor, \newblock Choosability in graphs, \newblock in \emph{Proceedings of the West Coast Conference on Combinatorics, Graph Theory and Computing}, Congressus Numerantium XXVI (1979), 125--157.
\bibitem[FKK2014]{FKK2014} Z.~F\"uredi, A.~Kostochka and M.~Kumbhat, \newblock Choosability with separation of complete multipartite graphs and hypergraphs, \newblock \emph{J. Graph Theory} 76 (2014), 129--137.
\bibitem[GuJiangWoodZhu2023]{GuJiangWoodZhu2023} Y.~Gu, Y.~Jiang, D.~R. Wood and X.~Zhu, \newblock Refined list version of Hadwiger's conjecture, \newblock \emph{SIAM J. Discrete Math.} 37 (2023), 1738--1750.
\bibitem[GuZhu2023]{GuZhu2023} Y.~Gu and X.~Zhu, \newblock Girth and $\lambda$-choosability of graphs, \newblock \emph{J. Graph Theory} 103 (2023), 493--501.
\bibitem[KTV1998Brooks]{KTV1998Brooks} J.~Kratochv{\'\i}l, Zs.~Tuza and M.~Voigt, \newblock Brooks-type theorems for choosability with separation, \newblock \emph{J. Graph Theory} 27 (1998), 43--49.
\bibitem[KTV1998Complexity]{KTV1998Complexity} J.~Kratochv{\'\i}l, Zs.~Tuza and M.~Voigt, \newblock Complexity of choosing subsets from color sets, \newblock \emph{Discrete Math.} 191 (1998), 139--148.
\bibitem[Kang2013]{Kang2013} R.~J. Kang, \newblock Improper choosability and Property B, \newblock \emph{J. Graph Theory} 73 (2013), 342--353.
\bibitem[KralSgall2005]{KralSgall2005} D.~Kr\'al' and J.~Sgall, \newblock Coloring graphs from lists with bounded size of their union, \newblock \emph{J. Graph Theory} 49 (2005), 177--186.
\bibitem[ST2015]{ST2015} D.~Saxton and A.~Thomason, \newblock Hypergraph containers, \newblock \emph{Invent. Math.} 201 (2015), 925--992.
\bibitem[Vizing1976]{Vizing1976} V.~G. Vizing, \newblock Coloring the vertices of a graph in prescribed colors, \newblock \emph{Diskret. Analiz} 29 (1976), 3--10 (in Russian).
\bibitem[Zhu2020]{Zhu2020} X.~Zhu, \newblock A refinement of choosability of graphs, \newblock \emph{J. Combin. Theory Ser. B} 141 (2020), 143--164.
\bibitem[ZhuZhu2021]{ZhuZhu2021} J.~Zhu and X.~Zhu, \newblock Chromatic $\lambda$-choosable and $\lambda$-paintable graphs, \newblock \emph{J. Graph Theory} 98 (2021), 642--652.
\bibitem[ZhuZhu2025]{ZhuZhu2025} J.~Zhu and X.~Zhu, \newblock Minimum non-chromatic-$\lambda$-choosable graphs, \newblock \emph{J. Graph Theory} 110 (2025), 283--289.
\end{thebibliography}
\end{document}
